% Source: https://artofproblemsolving.com/community/c3h1737089_00_what_is_the_answer_to_this
% Stable identity: aops-topic-1737089-q01
For every $x>0$ we have $x^0=1$, so $\lim_{x\to0^+}x^0=1$. For every $y>0$ we have $0^y=0$, so $\lim_{y\to0^+}0^y=0$.

If $f$ were continuous at $(0,0)$ with some value $f(0,0)=c$, then $f$ would tend to $c$ along every path into $(0,0)$. Along the path $(x,0)$, $x\to0^+$, the values are all $1$, forcing $c=1$; along the path $(0,y)$, $y\to0^+$, the values are all $0$, forcing $c=0$. This is impossible, so no choice of $0^0$ makes $x^y$ continuous at $(0,0)$; this is why $0^0$ is called an indeterminate form in analysis. (Other paths give other limits too: along $y=\dfrac{1}{\ln\dfrac{1}{x}}$ the value $x^y=e^{-1}$ is constant.)

In algebra and combinatorics, however, $a^0$ means the empty product, which is $1$ for every $a$, and with $0^0=1$ formulas such as $(1+x)^n=\sum_k\binom nk x^k$ at $x=0$, $\sum_k a_kx^k$ at $x=0$, and ``the number of functions from an $m$-set to an $n$-set is $n^m$'' (one function from the empty set) hold without exceptions. So $0^0=1$ is the useful convention there, while in limits $0^0$ has no single value.
